\conclusions{

    La presente investigación demostró la factibilidad, pertinencia y potencial 
    científico de aplicar la Ingeniería de Líneas de Productos (\ac{SPL}) a la 
    planificación y gestión curricular en educación superior. El análisis conceptual 
    permitió establecer una analogía formal entre los constructos fundamentales de 
    \ac{SPL} y las entidades curriculares, revelando que la variabilidad presente 
    en los planes de estudio —en forma de optatividades, prerrequisitos, menciones, 
    orientaciones y rutas formativas— puede ser correctamente modelada como un c
    onjunto estructurado y verificable mediante Modelos de Características (\ac{FM}). 
    Este hallazgo no solo valida el enfoque metodológico adoptado, sino que además 
    aporta un marco riguroso para la representación y análisis sistemático de la 
    complejidad curricular.

    Desde el punto de vista arquitectónico, la propuesta demostró que es viable 
    construir una plataforma computacional que haga uso de \ac{FM}s versionados, 
    analizados y validados de manera automática, con soporte para configuraciones 
    curriculares personalizadas. Para ello, se integraron técnicas de análisis 
    simbólico, resolución \ac{SAT}/\ac{SMT}, heurísticas evolutivas y modelado 
    estructural basado en grafos, garantizando precisión lógica y escalabilidad. 
    La introducción de un modelo híbrido de validación permite abordar escenarios 
    que combinan decisiones obligatorias, alternativas, exclusiones cruzadas y 
    restricciones globales, preservando coherencia y consistencia incluso ante 
    modificaciones evolutivas del currículo.

    La arquitectura resultante, fundamentada en patrones de diseño, desacoplamiento 
    conceptual y versionado inmutable, favorece la extensibilidad, la mantenibilidad y
    la auditabilidad. Asimismo, el almacenamiento estructurado de resultados 
    analíticos, métricas y configuraciones abre oportunidades para futuros sistemas 
    inteligentes de recomendación, análisis curricular comparativo e incluso 
    evaluación automatizada de rutas pedagógicas. Los resultados del modelado de 
    datos evidencian que la formalización curricular mediante \ac{FM}s no constituye 
    un ejercicio teórico aislado, sino una infraestructura aplicable en entornos 
    reales de gestión académica.

    Finalmente, la propuesta no pretende sustituir los procesos académicos 
    tradicionales, sino fortalecerlos mediante una abstracción formal que permita 
    control, evidencia verificable y exploración sistemática. Se concluye que la 
    integración de \ac{SPL} y currículo universitario representa un camino sólido 
    para la modernización de sistemas educativos, habilitando flexibilidad curricular 
    controlada, análisis automatizado, trazabilidad y personalización. Esta 
    contribución abre líneas futuras de investigación orientadas a (i) análisis
    semántico avanzado de restricciones, (ii) generación automática asistida de 
    itinerarios curriculares óptimos, (iii) interoperabilidad con sistemas 
    institucionales (SIS/\ac{LMS}) y (iv) incorporación de métricas pedagógicas 
    para evaluación dinámica de rutas de aprendizaje.

}